% =====================================================================
% GPU-Accelerated Graph-Colored Simulated Annealing for Integer Factorization
% First draft -- figures are textual placeholders (see \figplaceholder)
% =====================================================================
\documentclass[reprint,twocolumn,amsmath,amssymb,floatfix,nofootinbib]{revtex4-2}

\usepackage{graphicx}
\usepackage{amsmath}
\usepackage{amssymb}
\usepackage{booktabs}
% RevTeX 4-2 reimplements the float package, which breaks any package
% relying on \newfloat (such as the standard `algorithm' package).
% algorithm2e is self-contained and works under revtex.
\usepackage[ruled,vlined,linesnumbered]{algorithm2e}
\usepackage{xcolor}
\usepackage{hyperref}

% --- Placeholder figure boxes. First argument is the label, second is the
% --- caption / description. Replace with \includegraphics later.
\newcommand{\figplaceholder}[2]{%
  \begin{figure}[t]
    \centering
    \fbox{\parbox[c][3.2cm][c]{0.95\columnwidth}{\centering\itshape\small FIG: #2}}
    \caption{\textit{[placeholder]} #2}
    \label{#1}
  \end{figure}}
\newcommand{\figplaceholderwide}[2]{%
  \begin{figure*}[t]
    \centering
    \fbox{\parbox[c][4.5cm][c]{0.95\textwidth}{\centering\itshape\small FIG: #2}}
    \caption{\textit{[placeholder]} #2}
    \label{#1}
  \end{figure*}}

% --- Convenience macros ----------------------------------------------
\newcommand{\bN}{\mathbf{N}}
\newcommand{\dE}{\Delta E}
\newcommand{\bigO}{\mathcal{O}}

\begin{document}

\title{GPU-Accelerated Graph-Colored Simulated Annealing
for Integer Factorization}

\author{Advith Desu}
\author{Aryan Namboodri}
\author{Anil Prabhakar}
\affiliation{Centre for Quantum Information, Communication and Computing,\\
Indian Institute of Technology Madras, Chennai 600036, India}

\date{\today}

\begin{abstract}
The security of RSA rests on the hardness of factoring a semiprime
$N = PQ$. Shor's algorithm dissolves that hardness, but only at
qubit counts that remain decades away. In the interim, classical and
quantum-inspired Ising machines offer a near-term lever on the same
problem. We present an end-to-end pipeline that maps factorization to
a sparse Ising model, solves it on a single NVIDIA GH200 GPU with a
graph-colored simulated annealer, and closes the residual gap with a
guided brute-force search. The solver's three structural choices --- a
hub-tagged sparse CSR backed by a shared-memory hub cache, dense/sparse
spin binning matched to GPU resources, and Welsh--Powell coloring of the
interaction graph for exact parallel Metropolis --- together factor a
$\sim 100$-bit semiprime in $416$ seconds end-to-end, an
$8\times$ improvement over a $\sqrt{N}$-seeded brute-force baseline on
the same hardware.
\end{abstract}

\maketitle

% =====================================================================
\section{Introduction}\label{sec:intro}
% =====================================================================

Public-key cryptography assumes that some computational problems are
asymptotically hard. RSA in particular assumes that recovering the two
prime factors of a public modulus $N = PQ$ takes time exponential in
the bit length of $N$. Shor's algorithm collapses this assumption on a
fault-tolerant quantum computer, but the resource counts --- on the
order of millions of physical qubits for cryptographically relevant key
sizes~\cite{Jain2016} --- place that threat at least a decade
out. Meanwhile, classical and quantum-inspired Ising machines have
matured to the point where they routinely solve quadratic unconstrained
binary optimization (QUBO) problems with $10^4$--$10^5$ variables, and
recent work has factored integers as large as $2049$ bits on
specialized hardware~\cite{XiniGuo2026}.

Integer factorization can be cast as a QUBO~\cite{Jiang2018,Wang2020}.
The cleanest route is to write $N = PQ$ as a binary long
multiplication, treat the unknown bits of $P$, $Q$ and the carries as
binary variables, square the per-column constraints, and minimize the
resulting energy. The ground state encodes the factors. Whether that
ground state can actually be \emph{found} depends on three things: how
clean the resulting QUBO is, how well the solver exploits its sparsity,
and how the solver scales on the hardware it runs on.

This paper is about the second two. We treat the QUBO construction as
support infrastructure and focus on a GPU-resident simulated annealer
that has been shaped, end-to-end, by the structure of the Ising graphs
that factorization produces. Three design choices carry the
contribution:
\begin{enumerate}
  \item \textbf{Graph-colored parallel Metropolis.} A Welsh--Powell
    coloring of the interaction graph partitions spins into
    independent sets. Within a color, every spin's $\dE$ can be
    computed and applied in parallel without invalidating the global
    energy accumulator.
  \item \textbf{Dense/sparse spin binning.} Empirically, only a small
    set of "hub" spins (with row degree $\gtrsim 100$) attracts most
    of the coupling weight, while the rest of the graph stays sparse.
    The annealer routes the two populations through different kernels
    --- block-level shared-memory reduction for hubs, warp-level
    shuffle reduction for sparse spins --- sized to the GPU's actual
    resources.
  \item \textbf{Hub-tagged sparse CSR.} The high bit of each column
    index in the sparse-spin CSR encodes whether the neighbor is a hub
    (read from a shared-memory cache) or a non-hub (read from global
    memory). The hot path is a single branchless shared-memory load.
\end{enumerate}
Together, these let one GH200 sweep a $\sim 10^5$-spin Ising graph at
$\sim 0.1$\,ms per iteration, and factor a $\sim 100$-bit semiprime
end-to-end in $416$\,s.

Section~\ref{sec:formulation} casts factorization as a binary
optimization problem and works through a small example.
Section~\ref{sec:preproc} describes the classical simplification and
quadrization that turns it into a sparse Ising model.
Section~\ref{sec:scaling} shows how the structure of that model scales,
which motivates the kernel design in Section~\ref{sec:sa}.
Section~\ref{sec:bruteforce} describes the post-processing search that
closes the residual gap. Section~\ref{sec:results} reports timing and
solution-quality results across $16$--$128$-bit semiprimes, and
Section~\ref{sec:discussion} discusses limits and next steps.

% =====================================================================
\section{From factoring to a binary energy}\label{sec:formulation}
% =====================================================================

Let $N = PQ$ be a semiprime with $P, Q$ odd primes of equal bit length
$n = \lceil \tfrac{1}{2}\log_2 N \rceil$. Write
\begin{equation}
  P = \sum_{k=0}^{n-1} 2^k p_k,
  \qquad
  Q = \sum_{k=0}^{n-1} 2^k q_k,
  \qquad p_k,q_k \in \{0,1\}.
\end{equation}
Four bits are known a priori: $p_0 = q_0 = 1$ (both factors are odd)
and $p_{n-1} = q_{n-1} = 1$ (each is exactly $n$ bits).

The product $PQ$, expanded as a binary long multiplication, places the
cross-term $p_a q_b$ in column $a+b$. Carries propagate left; we
introduce binary carry variables $s_{i,j} \in \{0,1\}$, where $s_{i,j}$
is the bit of column $i$'s carry destined for column $j > i$ (with
weight $2^{j-i}$ when it arrives). Equating the value of column $i$ in
$PQ$ to the value of column $i$ in $N$ yields a per-column constraint:
\begin{equation}
  C_i \;=\; N_i \;-\; \sum_j q_j p_{i-j}
            \;-\; \sum_j s_{j,i}
            \;+\; \sum_j 2^{j} s_{i,i+j} \;=\; 0,
  \label{eq:Ci}
\end{equation}
where $N_i$ is the $i$th bit of $N$. The correct factor assignment
satisfies $C_i = 0$ for every column. Squaring and summing turns the
constraint system into an unconstrained optimization:
\begin{equation}
  E(\{p_i, q_i, s_{ij}\}) \;=\; \sum_i C_i^2.
  \label{eq:energy}
\end{equation}
The ground state $E = 0$ recovers the factors.

\paragraph{Worked example, $N = 391 = 17 \times 23$.}
With $n = 5$, the unknown factor bits are $p_1, p_2, p_3$ and
$q_1, q_2, q_3$ --- six free bits. The carry variables span at most
two columns ahead (each column sum is bounded by a small constant), so
the full system has nine column constraints $C_0,\dots,C_8$ and a
handful of carry variables. Three of those constraints are immediately
useful: $C_0$ is trivially satisfied ($1 - 1 = 0$); $C_8$ reduces to
$s_{6,8} + s_{7,8} = 0$, forcing both carries to zero; and $C_1$
reads $p_1 + q_1 = 1 + 2 s_{1,2}$, which has no solution with
$s_{1,2}=1$ (the LHS is at most $2$) so $s_{1,2} = 0$ and
$q_1 = 1 - p_1$. The same logic threads through $C_2$. Six rules of
this kind --- enumerated next --- handle the bulk of the
simplification.

% =====================================================================
\section{Classical pre-processing}\label{sec:preproc}
% =====================================================================

Squaring \eqref{eq:Ci} directly produces a polynomial with
\emph{many} more terms than needed: each carry variable shows up in
multiple clauses, and cross-products like $p_i q_j$ appear with very
different weights in different columns. Before handing the problem to
the annealer, we eliminate as many variables and as much redundancy as
possible using deterministic algebraic rules~\cite{Jiang2018}.

\subsection{Reduction rules}\label{sec:rules}

Each $C_i = 0$ is a linear (or low-degree) equation over $\{0,1\}$
variables. Six rules suffice to drive most of the carries and several
of the factor bits to fixed values; their formal statements and short
proofs are collected in Appendix~\ref{app:rules}.
\begin{description}
  \item[Saturated sum.] $\sum x_i = n$ (all coefficients equal $\pm 1$,
    $n$ terms) forces every $x_i = 1$.
  \item[Zero sum.] $\sum x_i = 0$ with same-sign coefficients forces
    every $x_i = 0$.
  \item[Doubling.] $x_1 + x_2 = 2 x_3$ forces $x_1 = x_2 = x_3$.
  \item[Coefficient bound.] If a single coefficient exceeds the
    achievable opposing sum, that term is forced to $0$ or $1$.
  \item[Complement.] $x_1 + x_2 = 1$ gives $x_2 = 1 - x_1$ and
    $x_1 x_2 = 0$, eliminating a variable and killing one quadratic
    term in any clause that referenced both.
  \item[Parity.] Both sides of a column clause must have the same
    parity; this pins the parity of the small factor-bit sum against
    the (always-even) carry-out terms.
\end{description}

The simplifier applies these rules in a fixed priority order, looping
until no rule fires on any clause. A seventh \emph{Replacement} rule
isolates an $s$-variable that appears with coefficient $\pm 1$ and
substitutes it back into every other clause. Replacement is powerful
but greedy: a long chain of cascading replacements can fuse many small
clauses into a single mega-clause whose squaring blows up the QUBO. We
guard against this with a \emph{checkpoint}: the simplifier saves its
state before the first Replacement and restores it if the chain failed
to produce new value assignments.

\subsection{Quadrization}\label{sec:quad}

The energy~\eqref{eq:energy} is degree four in the binary
variables. Cubic and quartic terms arise from squaring clauses that
contain products like $p_i q_j$. Standard
QUBO solvers (and Ising hardware) require degree two, so we apply
\emph{Rosenberg substitutions}~\cite{Rosenberg1975}: introduce an
auxiliary $w$ and add a penalty $P(A,B,w) = M(AB - 2Aw - 2Bw + 3w)$
with $M$ large enough that any minimizer has $w = AB$. After
substitution every cubic $AB \cdot s$ becomes a quadratic $w \cdot s$
plus a constant-size quadratic penalty. The two sign variants of the
cubic case, and the two variants of the quartic case (negative quartic
needs only a single auxiliary; positive quartic needs two and is
handled recursively), are standard. Implementation details and
coefficient bookkeeping follow~\cite{Jiang2018}.

\subsection{From QUBO to Ising}\label{sec:qubotoising}

A linear change of variables $\sigma_i = 1 - 2 x_i$ maps each binary
$x_i \in \{0,1\}$ to an Ising spin $\sigma_i \in \{-1, +1\}$. The QUBO
$H = \sum_i Q_{ii} x_i + \sum_{i < j} Q_{ij} x_i x_j$ becomes
\begin{equation}
  H(\boldsymbol\sigma) = \sum_i h_i \sigma_i
                       + \sum_{i<j} J_{ij} \sigma_i \sigma_j
                       + \text{const.},
\end{equation}
with, treating $Q$ symmetrically so $Q_{ij} = Q_{ji}$ for $i \neq j$,
\begin{equation}
  J_{ij} = \tfrac{1}{4} Q_{ij}, \quad
  h_i    = -\tfrac{1}{2} Q_{ii} - \tfrac{1}{4} \sum_{j \neq i} Q_{ij}.
\end{equation}
The Ising form is what every downstream piece consumes.

% =====================================================================
\section{Structural scaling and sparse storage}\label{sec:scaling}
% =====================================================================

What does the resulting Ising problem actually look like? Two
empirical facts shape the rest of the paper.

\paragraph{The coupling matrix is sparse, and stays sparse.}
\figplaceholderwide{fig:sparsity}{Sparsity of $J_{ij}$ at $N=39\,203$ (162 spins),
$159\,197$ (379), $338\,699$ (563), and $377\,977$ (902). Most
nonzeros concentrate in a small band along the top rows / left
columns, with a thin diagonal tail. Pattern persists across
orders of magnitude in $N$.}
Figure~\ref{fig:sparsity} plots the nonzero structure of $J_{ij}$
for four representative semiprimes. Two visual features dominate every
panel: a small dense block of "hub" rows / columns near index zero,
and a long thin diagonal trail. The dense block corresponds to
auxiliary $w$-variables introduced during quadrization
(Section~\ref{sec:quad}), each of which couples to every clause it
participates in; the diagonal trail comes from the carries
$s_{i,j}$, which couple only to their immediate column neighbors.

\paragraph{Hub spin count grows only with $\log N$.}
\figplaceholder{fig:hubcount}{Number of total Ising spins versus number of
\emph{densely connected} (row degree $\geq 128$) hub spins as the bit
width of $N$ grows from $16$ to $\sim 128$. Total grows as
$\bigO(n^3)$; hub count grows as $\bigO(\log N)$, plateauing near $120$
at $128$ bits.}
Figure~\ref{fig:hubcount} separates the total spin count from the
\emph{dense} spin count (row degree $\geq 128$). The former grows
roughly cubically with $n$; the latter grows only logarithmically, and
saturates around $120$ at $128$ bits. The dense set is small enough to
hold in shared memory in its entirety, and large enough that every
sparse spin reads from it on every sweep. This is the structural fact
that the kernel design in Section~\ref{sec:sa} exploits.

\paragraph{CSR storage.}
A dense $|J|^2$ representation is wasteful at $|J| > 10^4$ and
infeasible at $|J| > 10^5$. We store $J$ in Compressed Sparse Row
(CSR) format: a \texttt{row\_ptr} of $N+1$ row offsets, a
\texttt{col\_idx} array of column indices, and a \texttt{values} array
of the nonzeros. Memory is $\bigO(N + \mathrm{nnz})$ rather than
$\bigO(N^2)$; for the problems considered here, $\mathrm{nnz}$ is
typically $10$--$100\times$ smaller than $N^2$.

% =====================================================================
\section{Graph-colored simulated annealing on the GPU}\label{sec:sa}
% =====================================================================

The annealer is a Metropolis sampler that tries to flip every spin
once per sweep, accepting flips with probability
$\min(1, e^{-\beta \dE_i})$. The serial cost of one sweep is
$\bigO(\mathrm{nnz})$. To turn that into wall-clock time on a GPU we
need three things: a way to flip many spins in parallel without
breaking the energy bookkeeping, kernels whose work pattern matches
GPU resources, and a memory layout that keeps the hot data in fast
storage.

\subsection{Why parallel Metropolis needs coloring}\label{sec:sa-coloring}

Single-spin Metropolis updates the state one spin at a time. Two
spins coupled by $J_{ij} \neq 0$ cannot be flipped simultaneously: the
$\dE$ each one computes depends on the current value of the other,
so flipping both at once leaves the running energy accumulator
inconsistent with the new spin state.

\emph{Graph coloring} of the interaction graph (nodes = spins, edges =
nonzero $J_{ij}$) partitions the spins into independent sets such that
no two same-color spins are coupled. All spins of one color can then
be flipped in parallel: each one's $\dE$ depends only on spins of
\emph{other} colors, which are unchanged during this color's
pass. Sweeping every color in sequence gives a full Metropolis sweep
with no synchronization inside a color and one synchronization between
colors.

We compute the coloring once at setup using a Welsh--Powell greedy
heuristic (largest-degree-first, first-fit). The number of colors is
upper-bounded by $\Delta + 1$ where $\Delta$ is the maximum degree;
empirically the heuristic produces colorings within a few colors of
optimum.

\subsection{Dense and sparse spin bins}\label{sec:sa-bins}

The two empirical features from
Section~\ref{sec:scaling} --- a small dense hub set and a long sparse
tail --- map naturally onto two GPU kernels.
\begin{description}
  \item[Dense bin] (row degree $\geq 128$). One CUDA block of $1024$
    threads handles one spin. Each thread accumulates a partial sum
    over a stride of the spin's neighbor row; a shared-memory tree
    reduction collapses the partials to $\dE$. The block size matches
    the dense-row length so threads stay busy.
  \item[Sparse bin] (row degree $< 128$). One warp ($32$ threads)
    handles one spin; $32$ spins are packed into a $1024$-thread
    block. Each warp does a warp-shuffle reduction --- no shared
    memory, no explicit synchronization within the warp. This gives
    much higher occupancy than the dense layout would for short rows.
\end{description}
The threshold $128$ matches the warp's natural reduction tier and is
robust across the bit-width range we tested: above it, dense reduction
amortizes the block-launch overhead; below it, the warp layout wins on
occupancy.

\subsection{Hub-tagged sparse CSR}\label{sec:sa-hubs}

The dense (hub) spins are read by essentially every sparse spin in
every color pass. Reading them from global memory each time is
expensive; we cache them once per color pass.

Hub values for the entire problem are copied into a shared-memory
array (at most $256$ bytes total) at the start of each sparse kernel
launch. To resolve "is this neighbor a hub or not?" branchlessly, we
\emph{tag} the column-index array of the sparse-only CSR: the high
bit of \texttt{col\_idx[k]} encodes the answer.
\begin{itemize}
  \item High bit set: the low 31 bits are an index into the hub cache
    in shared memory.
  \item High bit clear: the entry is a raw global spin id, used as a
    fallback for the rare non-hub neighbor.
\end{itemize}
The inner neighbor-sum loop is therefore a sign test on
\texttt{col\_idx[k]}, a single shared-memory load on the hot path, and
a global-memory load on the cold path. Whenever a hub spin flips
(which can only happen on a color pass that contains hubs), the
hub-value cache is refreshed by a small dedicated kernel before any
sparse pass reads it.

\subsection{Auto start temperature and cooling}\label{sec:sa-temp}

A simulated-annealing run is parameterized by a temperature schedule
$T_t = T_0 \alpha^t$ with $0 < \alpha < 1$ and an inner sweep count.
The wrong $T_0$ wastes sweeps: too high and the run is effectively a
random walk for many iterations; too low and the system freezes into
the initial basin.

We estimate $T_0$ adaptively, in the style of
Ben-Ameur~\cite{BenAmeur2004}. Across $K = 10$ random spin
configurations we sample every spin's uphill $\dE$ (the energy cost of
a single flip when it raises the energy), average them, and set
\begin{equation}
  T_0 \;=\; \langle \dE^{+} \rangle \, / \, \ln(1 / \chi),
\end{equation}
where $\chi = 0.5$ is the target acceptance rate at the start of the
schedule. The sampling reuses the same GPU buffers and hub cache as
the annealing loop, so the cost is at most one full sweep per
configuration. Cooling is geometric with $\alpha = 0.95$; the schedule
runs until $T < 10^{-3}$.

\subsection{Numerical scaffolding}\label{sec:sa-scaffolding}

Two correctness checks and one bookkeeping invariant keep the run
honest.
\begin{description}
  \item[Symmetry check.] On load, the annealer samples random rows of
    the CSR and verifies that for every stored $(i, j, J_{ij})$ the
    partner $(j, i, J_{ji})$ is also stored and has the same value.
    The energy and $\dE$ formulas assume this; an upper-triangular-only
    CSR would silently halve every off-diagonal coupling.
  \item[Running-energy accumulator.] Instead of recomputing the total
    energy from scratch after each sweep, the annealer keeps a running
    sum: every accepted flip atomically adds its $\dE$ to a
    \texttt{double} accumulator. At the end of each annealing run, the
    accumulator is checked against a fresh full-state recomputation;
    relative drift stays below $10^{-9}$ for all problem sizes we
    tested.
  \item[Best-state mirror.] A second device-side spin buffer mirrors
    the live state whenever a new best running energy is observed.
    The output spin file is written from this mirror, so the reported
    state is the best one ever visited, not whatever happened to be
    live at the final $T$ step.
\end{description}

\subsection{Putting it together}\label{sec:sa-loop}

Algorithm~\ref{alg:setup} lists the one-time setup; the per-temperature
loop is Algorithm~\ref{alg:anneal}.

\begin{algorithm}[t]
  \caption{One-time setup before annealing}
  \label{alg:setup}
  Verify $J$ symmetry and zero diagonal (sampled)\;
  Bin each spin: \textbf{dense} if $d_i \geq 128$, else \textbf{sparse}\;
  $\mathrm{color}[\cdot] \leftarrow$ Welsh--Powell greedy
    coloring (largest-degree-first, first-fit)\;
  Bucket spins by color into dense / sparse offset arrays\;
  Build the hub-tagged sparse-only CSR\;
  Upload CSR, color tables, and hub arrays to the GPU\;
  Initialize a random spin state and compute the initial energy\;
  Estimate $T_0 = \langle \dE^{+} \rangle / \ln(1/\chi)$\;
\end{algorithm}

\begin{algorithm}[t]
  \caption{Color-parallel simulated annealing loop}
  \label{alg:anneal}
  \ForEach{$\beta$ in geometric schedule}{
    \For{$\mathrm{sweep} = 1 \dots m$}{
      Generate fresh uniform randoms for all $N$ spins\;
      \For{$c = 0 \dots N_{\mathrm{colors}} - 1$}{
        Launch \textsc{collectFlipCandidates}$\langle$dense$\rangle$ on color $c$\;
        Launch \textsc{collectFlipCandidates}$\langle$sparse$\rangle$ on color $c$\;
        \textsc{applyAllFlipsInColor}: flip every accepted
          spin, atomically add $\dE_i$ to the energy accumulator\;
        \If{color $c$ contained any hub}{Refresh the hub-value cache\;}
      }
    }
    \If{running energy $<$ best energy}{snapshot best spins\;}
  }
  \Return best-energy spin configuration\;
\end{algorithm}

% =====================================================================
\section{Brute-force post-processing}\label{sec:bruteforce}
% =====================================================================

The annealer rarely lands exactly on the ground state at the bit
widths we test. It does, however, land \emph{close}: post-processing
of its spin output recovers approximate factors $(P_{\mathrm{pred}},
Q_{\mathrm{pred}})$ that differ from the true $(P, Q)$ by a few
percent (Section~\ref{sec:results}). A targeted brute-force search
closes the residual gap.

\paragraph{Centering.}
The annealer provides two independent estimates of $P$:
$P_{\mathrm{pred}}$ from the spins directly, and
$\lfloor N / Q_{\mathrm{pred}} \rfloor$ from the complementary factor.
Their average is the search centroid:
\begin{equation}
  P_{\mathrm{centre}} \;=\;
  \tfrac{1}{2}\bigl(P_{\mathrm{pred}} + \lfloor N / Q_{\mathrm{pred}} \rfloor\bigr).
\end{equation}
Empirically, averaging halves the typical gap (see
Figure~\ref{fig:quality}).

\paragraph{Outward-spiral search.}
Candidates are tested in chunks of $C = 2^{14}$, in outward-spiral
order: chunk $k$ covers $[P_{\mathrm{centre}} + \tfrac{k}{2} C,
P_{\mathrm{centre}} + (\tfrac{k}{2}+1) C - 1]$ if $k$ is even (upward
step) and the analogous downward window if $k$ is odd. Threads claim
the next chunk via an atomic counter, so the search is partially
parallel and the closest candidates to $P_{\mathrm{centre}}$ are
always tested first.

\paragraph{Wheel sieve.}
Within a chunk, candidates that are divisible by $2$, $3$, $5$, $7$,
or $11$ are skipped without a divisibility test against $N$. The
wheel of modulus $2 \cdot 3 \cdot 5 \cdot 7 \cdot 11 = 2310$ has
$480$ residues coprime to all five small primes --- roughly $21\%$ of
$2310$ --- so the sieve filters out about $79\%$ of candidates at
essentially zero cost.

\paragraph{Early termination.}
The first thread to find $p \mid N$ sets a shared stop flag; all other
threads exit at the next chunk boundary. The expected wall-clock cost
follows the empirical model
\begin{equation}
  T_{\mathrm{bf}} \;\approx\;
  \frac{3.32 \times 10^{-9}}{\mathrm{threads}} \;\cdot\;
  \mathrm{Gap}(\%) \;\cdot\; 2^n,
  \label{eq:bf-time}
\end{equation}
where $n$ is the bit length of the (smaller) factor and
$\mathrm{Gap}(\%)$ is the relative distance between
$P_{\mathrm{centre}}$ and the true $P$. The exponential in $n$ is
unavoidable; the only knob the annealer controls is the multiplicative
prefactor.

% =====================================================================
\section{Results}\label{sec:results}
% =====================================================================

All experiments run on a single NVIDIA GH200 Grace Hopper Superchip
(Hopper architecture, $\mathrm{sm}\_90$). The annealing kernel is
compiled with \texttt{nvcc} (CUDA 12.x); QUBO construction and
brute-force post-processing run on the Grace CPU side. Cooling
schedule: geometric with $\alpha = 0.95$, $T_0$ auto-estimated at
$\chi = 0.5$, $T_{\min} = 10^{-3}$, $10$ sweeps per $\beta$ step.

For each bit width $n \in \{16, 20, 24, \dots, 124, 128\}$ we
generate a random $n$-bit semiprime $N = PQ$ with $P, Q$ random
$\lceil n/2 \rceil$-bit primes, build the QUBO once, and run the SA
ten times with independent seeds. Reported numbers use the winning
(lowest-energy) seed; mean and spread are reported where useful.

\subsection{Per-iteration throughput}

\figplaceholder{fig:throughput}{SA wall-clock per sweep iteration versus number of
Ising spins, comparing the CUDA kernel on GH200, a PyTorch port on
GH200, and the same PyTorch code on an Apple M2 Max via MPS. CUDA stays
near $0.05$\,ms across four decades of spin count; the PyTorch
variants drift upward past $10^3$ spins.}

Figure~\ref{fig:throughput} compares wall-clock per sweep across
the three implementations. The CUDA kernel stays near $0.05$\,ms per
iteration across four orders of magnitude in spin count; the PyTorch
port pays roughly $5\times$ that on the same GPU because its
gather/scatter primitives cannot exploit the hub cache. The Mac MPS
curve drifts upward earlier and faster, consistent with its lower
memory bandwidth.

\subsection{End-to-end construction + anneal time}

\figplaceholder{fig:scaling}{Time to build $J$ and time to run one full annealing
schedule, both as a function of $n$ in $[16, 128]$ bits. Construction
is sub-second up to $\sim 60$ bits and reaches $\sim 60$\,s at $128$
bits. Annealing dominates above $\sim 60$ bits, reaching $\sim 220$\,s
at $128$ bits.}
Figure~\ref{fig:scaling} plots both wall-clock components against
bit width. QUBO construction is dominated by the algebraic
simplification at small $n$ and by the squaring + quadrization at
large $n$; SA wall-clock scales roughly linearly with the spin count,
which itself grows as $\bigO(n^3)$.

\subsection{Solution quality}

\figplaceholder{fig:quality}{Percentage gap between predicted and true $P$ as a
function of $n$, for three estimators: $\sqrt{N}$ (uniform baseline),
$P_{\mathrm{pred}}$ from SA, and the averaged centroid
$\tfrac{1}{2}(P_{\mathrm{pred}} + N / Q_{\mathrm{pred}})$. Means across
the tested range: $10.33\%$ for $\sqrt{N}$, $4.23\%$ for raw SA, and
$2.98\%$ for the centroid.}
Figure~\ref{fig:quality} reports the relative gap between the
predicted and the true factor for three estimators:
\begin{enumerate}
  \item $\sqrt{N}$ --- the uniform "I know nothing about the factors"
    baseline. Mean gap $10.33\%$ over the tested range.
  \item $P_{\mathrm{pred}}$ from the SA spin state directly. Mean gap
    $4.23\%$.
  \item The averaged centroid
    $\tfrac{1}{2}(P_{\mathrm{pred}} + \lfloor N / Q_{\mathrm{pred}} \rfloor)$.
    Mean gap $2.98\%$.
\end{enumerate}
Two facts matter for the next step. The SA output is consistently
better than $\sqrt{N}$, but only by a small constant factor; and the
free averaging step gains as much again. Both contribute the same
multiplicative factor to the brute-force prefactor in
Eq.~\eqref{eq:bf-time}.

\subsection{Brute-force time as a function of gap}

\figplaceholder{fig:bruteforce}{Time taken to find the exact factor by the wheel-sieve
brute-force search, as a function of the raw gap between
$P_{\mathrm{centre}}$ and the true $P$, on a $72$-thread CPU. Below a
$\sim 10^5$ gap the search completes in milliseconds (sieve
dominated); above $\sim 10^7$ it follows the $\mathrm{Gap} \cdot 2^n$
prediction of Eq.~\eqref{eq:bf-time}.}

Figure~\ref{fig:bruteforce} validates the timing model of
Eq.~\eqref{eq:bf-time}. Below a gap of $\sim 10^5$ the search is
sieve-dominated and finishes in milliseconds; above $\sim 10^7$ it
scales linearly with the raw gap as predicted.

\subsection{$\sim 100$-bit end-to-end case study}

For
\[
  N \;=\; 416\,749\,328\,744\,899\,275\,664\,046\,210\,629,
\]
the SA-guided pipeline reports:
\begin{center}
\begin{tabular}{lr}
  \toprule
  Stage & Wall-clock \\
  \midrule
  Build $J$ matrix (CSR)              & $25.3$\,s \\
  Simulated annealing ($\alpha = 0.95$, auto $T_0$) & $122.8$\,s \\
  Brute-force from $P_{\mathrm{centre}} = 5.79 \times 10^{14}$ & $267.9$\,s \\
  \midrule
  \textbf{Total}                      & $\mathbf{416.0}$\,\textbf{s} \\
  \bottomrule
\end{tabular}
\end{center}
The naive $\sqrt{N}$-seeded baseline on the same $72$-thread CPU
takes $3323.8$\,s, i.e.\ the SA pre-step buys an $8\times$ speedup
end-to-end. The recovered factors are
$P = 573\,858\,364\,692\,161$ and $Q = 726\,223\,323\,360\,389$.

% =====================================================================
\section{Discussion}\label{sec:discussion}
% =====================================================================

\subsection{Where the time goes}

Across the bit range we tested, three regimes dominate
end-to-end wall-clock at different scales:
\begin{itemize}
  \item Below $\sim 40$ bits, the simplifier resolves the problem
    completely; no SA is required.
  \item Between $\sim 40$ and $\sim 80$ bits, SA dominates and the
    brute-force step is sub-second.
  \item Above $\sim 80$ bits, the multiplicative
    $\mathrm{Gap} \cdot 2^n$ factor of Eq.~\eqref{eq:bf-time} drives
    the brute-force step past the SA step.
\end{itemize}
The third regime is the operational bottleneck: SA hardware is
hardware-bounded, but the brute-force prefactor is purely a function
of solution quality. Cutting the gap in half halves the time at every
bit width.

\subsection{Limitations}\label{sec:limits}

Several caveats are worth being explicit about.

\paragraph{Solution-quality gap grows with $n$.}
The mean gap of $4.23\%$ (raw SA) is across $16$--$128$ bits; the
individual-instance gap is noisier at the upper end and we observed
runs with $> 8\%$ gap at $120$+ bits. The $\sim 100$-bit case study is
representative, not best-case.

\paragraph{$128$-bit hard cap.}
The QUBO construction code uses $\mathtt{\_\_uint128\_t}$ arithmetic
throughout. Going beyond $128$ bits requires switching to a multi-word
big-integer representation, with predictable but nontrivial code
changes.

\paragraph{Trial division is the wrong post-processing tool.}
At cryptographic key sizes (1024+ bits), a $\mathrm{Gap} \cdot 2^n$
brute-force cost is hopeless even with perfect predictions. The right
question --- which we leave to future work --- is what
factor-finding algorithm best exploits a $P_{\mathrm{centre}}$ within
a few percent of the truth. Pollard's $\rho$ and Lenstra's elliptic
curve method are both candidates that benefit much more from a good
starting point than linear scan does.

\paragraph{Single-instance scope.}
Each pipeline run targets one $N$. There is no notion of solving a
batch of independent semiprimes in one schedule, even though the GH200
has the memory for it. Batched annealing is a straightforward
extension we have not yet attempted.

\paragraph{This is not an RSA break.}
The largest semiprime factored here is $\sim 100$ bits; production RSA
moduli are $2048$ bits or larger. The point of the work is to measure
how a carefully designed classical annealer scales on commodity GPU
hardware, not to threaten deployed cryptography.

\subsection{Future work}

The bottleneck identified in Section~\ref{sec:limits} suggests several
directions. \emph{Better QUBO encodings}: the simplification in
Section~\ref{sec:preproc} is purely algebraic; symmetry-aware or
sieve-aware reductions might collapse far more of the problem before
SA sees it. \emph{Smarter post-processing}: Pollard's $\rho$, Lenstra
ECM, or even GNFS sieving seeded by $P_{\mathrm{centre}}$ would
exploit the annealer's output far more effectively than linear scan
does. \emph{Hardware co-design}: the kernel structure described in
Section~\ref{sec:sa} maps naturally onto coherent Ising machines and
FPGA-based annealers, both of which provide higher coupling-update
throughput than a GPU for the same energy budget.

% =====================================================================
\section{Conclusion}\label{sec:conclusion}
% =====================================================================

A GPU-resident simulated annealer, structured around the empirical
sparsity of factorization-derived Ising graphs, factors $\sim 100$-bit
semiprimes in minutes on a single GH200. The contribution is not a
break of any cryptographic system, but a measurement of how far
commodity classical hardware reaches against a problem usually framed
in quantum terms, and a concrete kernel design --- graph-colored
Metropolis with hub-tagged sparse CSR --- that other Ising-problem
domains can pick up directly.

% =====================================================================
\appendix
\section{Reduction-rule proofs}\label{app:rules}
% =====================================================================

For $x_i \in \{0,1\}$:

\paragraph{Saturated sum.}
If $\sum_{i=1}^n x_i = n$ then every $x_i = 1$, since $n$ binary
variables sum to $n$ only when all are $1$.

\paragraph{Zero sum.}
If $\sum_{i=1}^n x_i = 0$ with same-sign coefficients then every
$x_i = 0$, since a sum of non-negative binaries vanishes only when
every term does.

\paragraph{Doubling.}
$x_1 + x_2 = 2 x_3$ forces $x_1 = x_2 = x_3$: the LHS is in
$\{0,1,2\}$ and the RHS in $\{0,2\}$, so the LHS must be even; the
only way two binaries sum to an even number is both equal, and that
shared value matches $x_3$.

\paragraph{Coefficient bound.}
In a clause $\sum_i c_i x_i + C = 0$ with $S_+ = \sum_{c_i > 0} c_i$
and $S_- = \sum_{c_j < 0} |c_j|$: a positive coefficient $c_i$ with
$c_i > S_- - C$ forces $x_i = 0$ (otherwise the LHS cannot reach
zero), and $c_i > S_+ + C$ forces $x_i = 1$. The negative-coefficient
cases are symmetric.

\paragraph{Complement.}
$x_1 + x_2 = 1$ has only two solutions $(0,1)$ and $(1,0)$, so
$x_2 = 1 - x_1$ and $x_1 x_2 = 0$.

\paragraph{Parity.}
Both sides of any linear clause must agree mod $2$. In a column
clause, all carry terms have even coefficients, so the parity of the
small factor-bit sum is pinned by the parity of $N_i$. Two factor
bits with odd coefficients must therefore either be equal (if $N_i$ is
even) or complementary (if $N_i$ is odd).

\paragraph{Power rule (used by squaring).}
For any $x \in \{0,1\}$ and $m \geq 1$, $x^m = x$. This makes the
squared energy expressible as a degree-2 QUBO once cubic and quartic
terms are quadrized.

% =====================================================================
\bibliography{paper}
% =====================================================================

\end{document}
